\documentclass{article}
\usepackage{amsmath,amsthm,amssymb,mathtools,fullpage}
\usepackage[colorlinks=true,linkcolor=blue,citecolor=blue]{hyperref}

%--------------------------------------------------------------------
% Theorem environments
%--------------------------------------------------------------------
\newtheorem{theorem}{Theorem}[section]
\newtheorem{lemma}[theorem]{Lemma}
\newtheorem{proposition}[theorem]{Proposition}
\newtheorem{corollary}[theorem]{Corollary}
\theoremstyle{definition}
\newtheorem{definition}[theorem]{Definition}
\theoremstyle{remark}
\newtheorem{remark}[theorem]{Remark}

%--------------------------------------------------------------------
% Macros
%--------------------------------------------------------------------
\newcommand{\R}{\mathbb{R}}
\newcommand{\N}{\mathbb{N}}
\newcommand{\e}{\varepsilon}
\newcommand{\half}{\tfrac{1}{2}}
\newcommand{\Lc}{\mathcal{L}}
\newcommand{\norm}[1]{\left\|#1\right\|}
\newcommand{\abs}[1]{\left|#1\right|}
\newcommand{\inner}[2]{\langle #1,\, #2\rangle}
\newcommand{\ddt}{\tfrac{d}{dt}}

\title{Phi-Renormalization for Axisymmetric-with-Swirl\\
Navier--Stokes: A Conditional Reduction of Global Regularity}
\author{Jonathan Robert Simons, CRNA \\
Prime Field Technologies LLC, Savannah, Georgia \\
\texttt{jonathansimons357@proton.me}}
\date{April 2026 \quad --- \quad Preprint}

\begin{document}

\maketitle

%====================================================================
\begin{abstract}
%====================================================================
We study the three-dimensional incompressible Navier--Stokes equations
with axisymmetric-with-swirl initial data on $\R^3$, augmented by the
linear fractional hyperdissipation $Q_1 = -\e(-\Delta)^{(\beta+2)/2}$
with $\beta = 0.6 > \frac{1}{2}$.  By the Lions threshold \cite{Lions1969},
the augmented system admits a unique global smooth solution for every
fixed $\e > 0$.

The central technical contribution is the \emph{Phi-renormalization}
(Theorem~\ref{thm:phi}):  the substitution $\Phi = u^\theta/r$, together
with $\Gamma = r u^\theta$, yields the identity
$r^{-4}\partial_z(\Gamma^2) = \partial_z(\Phi^2)$.  This rewrites the
swirl-vorticity source term without the explicit $1/r^4$ prefactor,
relocating the axis difficulty into the advection coefficient $u^r/r$;
it does not by itself remove that difficulty.

For the augmented system we establish: (i)~uniform-in-$\e$ $L^2$
energy and $\Gamma$-bounds; (ii)~a Phi-energy estimate for the renormalized
swirl $\Phi$ with an explicit Gronwall constant $C(\e)$ depending on
$\norm{u^r_\e/r}_{L^\infty}$; (iii)~weak compactness via Aubin--Lions as
$\e\to 0$; (iv)~every subsequential limit is a Leray--Hopf weak solution
of the classical system.  Whether this weak solution is globally smooth
reduces to a single explicitly stated open condition:
$\e$-independence of $C(\e)$ (Open Problem~\ref{op:gronwall}).

\medskip\noindent
\textbf{What this paper does not prove.}\enspace
Unconditional global regularity of the axisymmetric-with-swirl system: the
reduction is conditional on the $\e$-independent bound of Open
Problem~\ref{op:gronwall}.

\medskip\noindent
\textbf{MSC 2020:} 35Q30, 76D05, 76D03, 35B65.\\
\textbf{Keywords:} Navier--Stokes, axisymmetric swirl,
Phi-renormalization, axis singularity, fractional dissipation,
global regularity.
\end{abstract}

\tableofcontents

%====================================================================
\section{Introduction}
%====================================================================

The global regularity of the three-dimensional incompressible Navier--Stokes
equations
\begin{equation}\label{eq:NS}
  \partial_t u + (u\cdot\nabla)u = -\nabla p + \nu\Delta u,
  \quad \operatorname{div} u = 0,
  \quad u(0) = u_0,
\end{equation}
is one of the seven Clay Millennium Prize Problems \cite{FeffermanClay}.
Leray \cite{Leray1934} constructed global weak solutions; their smoothness
for all time remains unresolved.

In the \emph{axisymmetric-with-swirl} class, the flow has an additional
algebraic structure in cylindrical coordinates: the velocity field is
independent of the azimuthal angle $\theta$ and carries a nonzero swirl
component $u^\theta$.  The swirl introduces a centrifugal forcing term in
the vorticity equation that is potentially singular on the symmetry axis
$r = 0$.  Controlling this term is the central difficulty of the
axisymmetric-with-swirl problem; in the absence of swirl (i.e., $u^\theta
\equiv 0$), global regularity is known \cite{Ladyzhenskaya,UY1968}.

We study the $Q_1$-augmented system
\begin{equation}\label{eq:NSe}
  \partial_t u + (u\cdot\nabla)u = -\nabla p + \nu\Delta u
  - \e(-\Delta)^{(\beta+2)/2}u,
  \quad \operatorname{div} u = 0,
  \quad u(0) = u_0,
\end{equation}
with $\e > 0$ and $\beta = 0.6$.  The augmentation is a \emph{linear}
fractional hyperdissipation; with $\beta > \frac{1}{2}$, Lions
\cite{Lions1969} guarantees global smooth solutions for each fixed $\e > 0$.
Our contribution is to analyze the limit $\e\to 0$ for axisymmetric-with-swirl
data by exploiting the Phi-renormalization.

\subsection*{Organization}

Section~\ref{sec:notation} fixes notation.
Section~\ref{sec:axi} derives the axisymmetric-with-swirl scalar equations.
Section~\ref{sec:Q1} records the Lions theorem.
Section~\ref{sec:energy} establishes uniform $L^2$ and $L^\infty$ bounds.
Section~\ref{sec:phi} contains the Phi-renormalization.
Section~\ref{sec:phienergy} derives the Phi-energy estimate in full.
Section~\ref{sec:global} establishes global regularity of the $\e$-system.
Section~\ref{sec:compact} proves the Aubin--Lions compactness and the
limit theorem.
Section~\ref{sec:open} states the open problems precisely.
Section~\ref{sec:conclusion} concludes.

%====================================================================
\section{Notation and Function Spaces}\label{sec:notation}
%====================================================================

We work on $\R^3$ with cylindrical coordinates $(r,\theta,z)$ where
$r = \sqrt{x_1^2+x_2^2}\ge 0$.

\medskip\noindent
\textbf{Sobolev spaces.}
For $s\ge 0$, $H^s(\R^3)$ is the standard Sobolev space with
\[
  \norm{f}_{H^s}^2 = \int_{\R^3}(1+|\xi|^2)^s|\hat f(\xi)|^2\,d\xi,
  \qquad
  \norm{f}_{\dot H^s}^2 = \int_{\R^3}|\xi|^{2s}|\hat f(\xi)|^2\,d\xi.
\]

\medskip\noindent
\textbf{Weighted $L^2$ spaces.}
We use the measure $r\,dr\,dz$ throughout; explicitly,
\[
  \norm{f}_{L^2_r}^2 = \int_0^\infty\int_{-\infty}^\infty |f(r,z)|^2\,
  r\,dr\,dz
\]
and $H^1_r$ is defined with $\norm{\nabla f}_{L^2_r}$ as the gradient part
of the norm.  We also use the measure $dr\,dz$ (no weight) and write
$L^2(dr\,dz)$ to distinguish.

\medskip\noindent
\textbf{Axisymmetric-with-swirl fields.}
A divergence-free velocity field is axisymmetric-with-swirl if
\[
  u = u^r(r,z,t)\,\mathbf{e}_r + u^\theta(r,z,t)\,\mathbf{e}_\theta
    + u^z(r,z,t)\,\mathbf{e}_z.
\]
\emph{Axis regularity conditions:}  $u^r = u^\theta = \partial_r u^z = 0$
on $r = 0$.  These are compatible conditions preserved by the flow.

\medskip\noindent
\textbf{Key variables.}
\begin{align*}
  \Gamma &:= r\,u^\theta  && \text{(angular momentum per unit mass)},\\
  \Phi   &:= u^\theta/r   && \text{(renormalized swirl)},\\
  \omega^\theta &:= \partial_z u^r - \partial_r u^z
                          && \text{(azimuthal vorticity)},\\
  \Omega &:= \omega^\theta/r
                          && \text{(renormalized azimuthal vorticity)}.
\end{align*}
Note $\Gamma = r^2\Phi$ (equivalently $u^\theta = r\Phi$), and $\Phi$ extends
smoothly across $r=0$ because $u^\theta = O(r)$ at the axis.

\medskip\noindent
\textbf{Constants.}  $C$ denotes a positive constant depending only on
$\nu$, $\norm{u_0}_{L^2}$, and universal Sobolev constants.  Dependence
on $\e$ is always written $C(\e)$ and tracked explicitly.

%====================================================================
\section{The Axisymmetric-with-Swirl Equations}\label{sec:axi}
%====================================================================

\subsection{Scalar form in cylindrical coordinates}

For an axisymmetric-with-swirl velocity field the incompressibility
condition reduces to
\begin{equation}\label{eq:incomp}
  \frac{1}{r}\partial_r(r\,u^r) + \partial_z u^z = 0,
  \qquad\text{equivalently}\qquad
  \partial_r(r\,u^r) + \partial_z(r\,u^z) = 0.
\end{equation}

The $\theta$-component of the Navier--Stokes momentum equation is
\begin{equation}\label{eq:utheta}
  \partial_t u^\theta + u^r\partial_r u^\theta + u^z\partial_z u^\theta
  + \frac{u^r u^\theta}{r}
  = \nu\!\left(\partial_r^2 u^\theta + \frac{1}{r}\partial_r u^\theta
    - \frac{u^\theta}{r^2} + \partial_z^2 u^\theta\right).
\end{equation}
The centripetal term $u^r u^\theta/r$ on the left couples the swirl and
radial motion.

\subsection{The \texorpdfstring{$\Gamma$}{Gamma} equation}

\begin{lemma}\label{lem:Gammaeq}
  The angular momentum $\Gamma = r u^\theta$ satisfies
  \begin{equation}\label{eq:Gammaeq}
    (\partial_t + u^r\partial_r + u^z\partial_z)\Gamma
    = \nu\!\left(\partial_r^2\Gamma - \frac{1}{r}\partial_r\Gamma
      + \partial_z^2\Gamma\right).
  \end{equation}
\end{lemma}
\begin{proof}
Multiply \eqref{eq:utheta} by $r$.  The left side gives
\[
  r(\partial_t u^\theta + u^r\partial_r u^\theta + u^z\partial_z u^\theta)
  + u^r u^\theta
  = (\partial_t + u^r\partial_r + u^z\partial_z)(ru^\theta)
  = \frac{D\Gamma}{Dt},
\]
where $D/Dt = \partial_t + u^r\partial_r + u^z\partial_z$.  For the right
side, set $\Gamma = r u^\theta$ and compute using $\partial_r\Gamma = u^\theta
+ r\partial_r u^\theta$, $\partial_{rr}\Gamma = 2\partial_r u^\theta
+ r\partial_{rr}u^\theta$:
The $\theta$-component of the \emph{vector} Laplacian is
$\Delta_{\rm cyl}u^\theta - u^\theta/r^2$, where
$\Delta_{\rm cyl} = \partial_{rr} + \frac{1}{r}\partial_r + \partial_{zz}$
is the scalar Laplacian acting on the $\theta$-independent profile
$u^\theta$ and $-u^\theta/r^2$ is the single curvature term.  Thus
\begin{align*}
  \nu r\!\left(\Delta_{\rm cyl}u^\theta - \frac{u^\theta}{r^2}\right)
  &= \nu r\!\left(\partial_{rr}u^\theta + \frac{1}{r}\partial_r u^\theta
    + \partial_{zz}u^\theta - \frac{u^\theta}{r^2}\right)\\
  &= \nu\!\left(r\partial_{rr}u^\theta + \partial_r u^\theta
    + r\partial_{zz}u^\theta - \frac{u^\theta}{r}\right).
\end{align*}
Substituting $r\partial_{rr}u^\theta = \partial_{rr}\Gamma
- 2\partial_r u^\theta$, $\partial_r u^\theta = \partial_r\Gamma/r - \Gamma/r^2$,
and $u^\theta = \Gamma/r$:
\begin{align*}
  &= \nu\!\left(\partial_{rr}\Gamma - 2\partial_r u^\theta
    + \partial_r u^\theta - \frac{u^\theta}{r} + \partial_{zz}\Gamma\right)
  = \nu\!\left(\partial_{rr}\Gamma - \partial_r u^\theta
    - \frac{u^\theta}{r} + \partial_{zz}\Gamma\right)\\
  &= \nu\!\left(\partial_{rr}\Gamma
    - \Bigl(\tfrac{\partial_r\Gamma}{r} - \tfrac{\Gamma}{r^2}\Bigr)
    - \frac{\Gamma}{r^2} + \partial_{zz}\Gamma\right)
  = \nu\!\left(\partial_{rr}\Gamma - \frac{1}{r}\partial_r\Gamma
    + \partial_{zz}\Gamma\right).
\end{align*}
The two curvature contributions cancel exactly,
$+\Gamma/r^2 - \Gamma/r^2 = 0$, leaving no $\Gamma/r^2$ term and yielding
\eqref{eq:Gammaeq}. \quad (See also \cite{MajdaBertozzi}, Ch.~2.)
\end{proof}

\begin{remark}
  The operator on the right of \eqref{eq:Gammaeq} is
  $\Lc_0 = \partial_{rr} - (1/r)\partial_r + \partial_{zz}$.
  In the weighted space $L^2_r$, integration by parts gives
  \[
    \int\Lc_0\Gamma\cdot\Gamma\cdot r\,dr\,dz
    = -\norm{\nabla\Gamma}_{L^2_r}^2,
  \]
  so $\Lc_0$ is dissipative in $L^2_r$.  The $-(1/r)\partial_r$ term is a
  drift pointing away from the axis, which helps maintain coercivity.
\end{remark}

\subsection{The vorticity equation for \texorpdfstring{$\Omega$}{Omega}}

\begin{lemma}\label{lem:Omegaeq}
  The renormalized azimuthal vorticity $\Omega = \omega^\theta/r$ satisfies
  \begin{equation}\label{eq:Omegaeq}
    (\partial_t + u^r\partial_r + u^z\partial_z)\Omega
    = \frac{1}{r^4}\partial_z(\Gamma^2)
    + \nu\!\left(\partial_r^2\Omega + \frac{3}{r}\partial_r\Omega
      + \partial_z^2\Omega\right).
  \end{equation}
\end{lemma}
\begin{proof}
  The azimuthal component of the NS vorticity equation
  $\partial_t\omega + (u\cdot\nabla)\omega = (\omega\cdot\nabla)u
  + \nu\Delta_{\rm vec}\omega$ reduces for axisymmetric fields to
  \[
    (\partial_t + u^r\partial_r + u^z\partial_z)\omega^\theta
    - \frac{u^r\omega^\theta}{r}
    = \frac{\partial_z(u^\theta)^2}{r}
    + \nu\!\left(\Delta\omega^\theta - \frac{\omega^\theta}{r^2}\right).
  \]
  Divide by $r$ and write $\omega^\theta = r\Omega$.  The left side becomes
  $(\partial_t + u^r\partial_r + u^z\partial_z)\Omega$ (the $u^r/r$ terms
  cancel).  The source term: $(u^\theta)^2 = \Gamma^2/r^2$, so
  $\partial_z(u^\theta)^2/r^2 = \partial_z(\Gamma^2/r^2)/r^2
  = (1/r^4)\partial_z(\Gamma^2)$ (since $\partial_z$ holds $r$ fixed).
  
  For the diffusion term, compute $\nu r^{-1}(\Delta(r\Omega) - \Omega)$:
  \begin{align*}
    \Delta(r\Omega) &= \partial_{rr}(r\Omega) + \tfrac{1}{r}\partial_r(r\Omega)
      + \partial_{zz}(r\Omega)\\
    &= (2\partial_r\Omega + r\partial_{rr}\Omega)
      + (\tfrac{\Omega}{r} + \partial_r\Omega)
      + r\partial_{zz}\Omega.
  \end{align*}
  Thus $\Delta(r\Omega) - \Omega = r\partial_{rr}\Omega + 3\partial_r\Omega
  + \Omega/r + r\partial_{zz}\Omega - \Omega$.  Since $\Omega = O(r^0)$ at
  the axis, the $\Omega/r$ and $\Omega$ terms combine as
  $(1/r - 1)\Omega \to 0$; more precisely, dividing by $r$:
  \[
    \frac{1}{r}\bigl(\Delta(r\Omega)-\Omega\bigr)
    = \partial_{rr}\Omega + \frac{3}{r}\partial_r\Omega
      + \partial_{zz}\Omega,
  \]
  yielding \eqref{eq:Omegaeq}.
\end{proof}

\subsection{The equation for \texorpdfstring{$\Phi$}{Phi}}\label{subsec:Phieq}

\begin{lemma}\label{lem:Phieq}
  The renormalized swirl $\Phi = u^\theta/r$ satisfies
  \begin{equation}\label{eq:Phieq}
    \frac{D\Phi}{Dt} + \frac{2u^r\Phi}{r}
    = \nu\!\left(\partial_{rr}\Phi + \frac{3}{r}\partial_r\Phi
      + \partial_{zz}\Phi\right),
  \end{equation}
  where $D/Dt = \partial_t + u^r\partial_r + u^z\partial_z$.
\end{lemma}
\begin{proof}
  Write $\Gamma = r^2\Phi$.  Then
  \[
    \frac{D\Gamma}{Dt}
    = r^2\frac{D\Phi}{Dt}
      + u^r\partial_r(r^2)\Phi
    = r^2\frac{D\Phi}{Dt} + 2r u^r\Phi.
  \]
  Substituting $\Gamma = r^2\Phi$ into \eqref{eq:Gammaeq}:
  \begin{align*}
    \partial_{rr}(r^2\Phi) &= 2\Phi + 4r\partial_r\Phi + r^2\partial_{rr}\Phi,\\
    \frac{1}{r}\partial_r(r^2\Phi) &= 2\Phi + r\partial_r\Phi,\\
    \partial_{zz}(r^2\Phi) &= r^2\partial_{zz}\Phi.
  \end{align*}
  So $\nu\Lc_0\Gamma = \nu(2\Phi + 4r\partial_r\Phi + r^2\partial_{rr}\Phi
  - 2\Phi - r\partial_r\Phi + r^2\partial_{zz}\Phi)
  = \nu r^2(\partial_{rr}\Phi + \frac{3}{r}\partial_r\Phi
  + \partial_{zz}\Phi)$, using that \eqref{eq:Gammaeq} carries no
  $\Gamma/r^2$ term.  Setting
  $r^2\frac{D\Phi}{Dt} + 2ru^r\Phi = \nu r^2(\partial_{rr}\Phi
  + (3/r)\partial_r\Phi + \partial_{zz}\Phi)$ and dividing by $r^2$
  gives \eqref{eq:Phieq}.
\end{proof}

\subsection{The augmented scalar system}

Including the $Q_1$ fractional dissipation, the augmented equations for
$(\Omega_\e, \Gamma_\e, \Phi_\e)$ are:
\begin{align}
  (\partial_t + u_\e\cdot\nabla)\Omega_\e
    &= \frac{1}{r^4}\partial_z(\Gamma_\e^2)
       + \nu\Lc_4\Omega_\e
       - \e(-\Delta)^{1.3}\Omega_\e,
       \label{eq:NSe-Omega}\\
  (\partial_t + u_\e\cdot\nabla)\Gamma_\e
    &= \nu\Lc_0\Gamma_\e - \e(-\Delta)^{1.3}\Gamma_\e,
       \label{eq:NSe-Gamma}\\
  \frac{D\Phi_\e}{Dt} + \frac{2u_\e^r\Phi_\e}{r}
    &= \nu\Lc_4\Phi_\e - \e(-\Delta)^{1.3}\Phi_\e,
       \label{eq:NSe-Phi}
\end{align}
where $\Lc_4 = \partial_{rr} + (3/r)\partial_r + \partial_{zz}$ is the same
operator that governs $\Omega$ in \eqref{eq:Omegaeq}.

%====================================================================
\section{The \texorpdfstring{$Q_1$}{Q1} Augmentation
  and the Lions Threshold}\label{sec:Q1}
%====================================================================

\begin{theorem}[Lions 1969, \cite{Lions1969}]\label{thm:lions}
  Let $\beta \ge \tfrac{1}{2}$.  For each $\e > 0$ and
  axisymmetric-with-swirl $u_0\in H^s_{\rm axi}(\R^3)$ with $s$
  sufficiently large, the $Q_1$-augmented system \eqref{eq:NSe} admits a
  unique global smooth solution $u_\e \in C^\infty(\R^3\times[0,\infty))$.
\end{theorem}

Our value $\beta = 0.6$ satisfies $\beta > \tfrac{1}{2}$ strictly.
Theorem~\ref{thm:lions} is applied as a black box; the contribution of
this paper is the analysis of the limit $\e\to 0$ for axisymmetric-with-swirl
data via the Phi-renormalization.

%====================================================================
\section{Uniform Energy Estimates}\label{sec:energy}
%====================================================================

\subsection{\texorpdfstring{$L^2$}{L2} energy inequality}

\begin{proposition}\label{prop:L2}
  The $Q_1$-augmented solutions satisfy
  \begin{equation}\label{eq:energy}
    \half\ddt\norm{u_\e}_{L^2}^2
    + \nu\norm{\nabla u_\e}_{L^2}^2
    + \e\norm{u_\e}_{\dot H^{2.6}}^2 = 0.
  \end{equation}
  Consequently, $\norm{u_\e(t)}_{L^2} \le \norm{u_0}_{L^2} =: M$
  for all $t\ge 0$ and all $\e > 0$, and
  \begin{equation}\label{eq:H1int}
    \int_0^T\norm{\nabla u_\e}_{L^2}^2\,dt \le \frac{M^2}{2\nu},
    \qquad
    \e\int_0^T\norm{u_\e}_{\dot H^{2.6}}^2\,dt \le \frac{M^2}{2},
  \end{equation}
  both uniform in $\e > 0$.
\end{proposition}
\begin{proof}
  Take the $L^2(\R^3)$ inner product of \eqref{eq:NSe} with $u_\e$.
  The pressure term vanishes by $\operatorname{div}u_\e = 0$.  The
  nonlinear term vanishes:
  $\inner{(u_\e\cdot\nabla)u_\e}{u_\e} = \half\int u_\e\cdot\nabla|u_\e|^2
  = -\half\int(\operatorname{div}u_\e)|u_\e|^2 = 0$.
  The viscous term gives $-\nu\norm{\nabla u_\e}_{L^2}^2$, and the
  $Q_1$ term gives $-\e\norm{u_\e}_{\dot H^{2.6}}^2 \le 0$.
  This yields \eqref{eq:energy}.  Integrating in $t$ and using
  non-negativity of all terms gives \eqref{eq:H1int}.
\end{proof}

\subsection{\texorpdfstring{$L^\infty$}{L-infty} bound on
  \texorpdfstring{$\Gamma$}{Gamma}}

\begin{lemma}\label{lem:Gamma}
  If $\Gamma_0 = r\,u_0^\theta \in L^\infty\cap L^2(\R^3)$, then
  \begin{equation}\label{eq:Gammainf}
    \norm{\Gamma_\e(t)}_{L^\infty} \le \norm{\Gamma_0}_{L^\infty}
    =: G_0
    \quad\text{for all } t\ge 0,\ \e > 0.
  \end{equation}
\end{lemma}
\begin{proof}
  Equation \eqref{eq:NSe-Gamma} is
  $D\Gamma_\e/Dt = \nu\Lc_0\Gamma_\e - \e(-\Delta)^{1.3}\Gamma_\e$.
  Both $\nu\Lc_0$ and $-\e(-\Delta)^{1.3}$ are dissipative operators.
  Let $\Gamma_\e$ achieve its spatial maximum at $(r_0,z_0,t_0)$ with
  $r_0 > 0$.  At that point $\nabla\Gamma_\e = 0$ and
  $\partial_{rr}\Gamma_\e \le 0$, $\partial_{zz}\Gamma_\e \le 0$,
  giving $\Lc_0\Gamma_\e \le 0$.  The nonlocal term
  $-\e(-\Delta)^{1.3}\Gamma_\e \le 0$ at a global maximum by the
  nonlocal maximum principle \cite{Stinga2010}.  Hence $D\Gamma_\e/Dt \le 0$
  at the maximum, so the maximum is non-increasing.  On the axis $r=0$,
  $\Gamma_\e = r u_\e^\theta = 0$ by the axis regularity condition.
  Therefore $\norm{\Gamma_\e(t)}_{L^\infty}$ is non-increasing in $t$,
  giving \eqref{eq:Gammainf}.
\end{proof}

%====================================================================
\section{The Phi-Renormalization}\label{sec:phi}
%====================================================================

\subsection{The algebraic cancellation identity}

\begin{lemma}[Exact algebraic cancellation]\label{lem:cancel}
  With $\Gamma = r u^\theta$ and $\Phi = u^\theta/r$, so that
  $\Gamma = r^2\Phi$, the following identity holds for all $(r,z,t)$:
  \begin{equation}\label{eq:cancel}
    \frac{1}{r^4}\partial_z(\Gamma^2) = \partial_z(\Phi^2).
  \end{equation}
  In particular, the right-hand side is smooth on the axis $r=0$,
  since $\Phi = u^\theta/r$ extends smoothly there (the axis condition
  $u^\theta = O(r)$ is preserved by the flow).
  No Hardy inequality is required or used.
\end{lemma}
\begin{proof}
  Direct computation:  $\Gamma^2 = (r^2\Phi)^2 = r^4\Phi^2$.
  Differentiating in $z$ with $r$ held fixed:
  $\partial_z(\Gamma^2) = r^4\,\partial_z(\Phi^2)$.
  Dividing both sides by $r^4$ gives \eqref{eq:cancel}.
\end{proof}

\begin{remark}
  Identity \eqref{eq:cancel} is purely algebraic and holds
  independently of whether the flow is viscous, inviscid, or augmented.
  It is the sole mechanism behind the cancellation; no functional
  inequality is invoked.
\end{remark}

\subsection{Effect on the vorticity equation}

Substituting \eqref{eq:cancel} into \eqref{eq:NSe-Omega}, the augmented
$\Omega$-equation becomes
\begin{equation}\label{eq:Omega-renorm}
  (\partial_t + u_\e\cdot\nabla)\Omega_\e
  = \partial_z(\Phi_\e^2)
  + \nu\Lc_4\Omega_\e
  - \e(-\Delta)^{1.3}\Omega_\e.
\end{equation}
The source term is now $\partial_z(\Phi_\e^2) = 2\Phi_\e\partial_z\Phi_\e$,
which is smooth across $r=0$.

\begin{theorem}[Phi-renormalization]\label{thm:phi}
  For axisymmetric-with-swirl $Q_1$-augmented solutions with
  $\Gamma_0\in L^\infty\cap L^2$ and $\Phi_0 = u_0^\theta/r \in L^2_r$:
  \begin{enumerate}
  \item[\textup{(i)}]
    \emph{(No axis singularity.)}
    The singular term $r^{-4}\partial_z(\Gamma_\e^2)$ in the
    $\Omega$-equation is identically equal to $\partial_z(\Phi_\e^2)$
    by Lemma~\ref{lem:cancel}.  No residual singularity at $r=0$
    remains.

  \item[\textup{(ii)}]
    \emph{(Phi-energy estimate.)}
    There exists a constant $C(\e)$, bounded for each fixed $\e > 0$
    and depending on $\norm{u^r_\e/r}_{L^\infty_{t,x}}$, $G_0$, and
    $M$, such that
    \begin{equation}\label{eq:phi-energy}
      \half\ddt\norm{\Phi_\e}_{L^2_r}^2
      + \nu\norm{\nabla\Phi_\e}_{L^2_r}^2
      + \nu\int_{\R}\!\Phi_\e(0,z)^2\,dz
      + \e\norm{\Phi_\e}_{\dot H^{2.6}_r}^2
      \le C(\e)\norm{\Phi_\e}_{L^2_r}^2.
    \end{equation}

  \item[\textup{(iii)}]
    \emph{(Gronwall bound.)}
    By Gronwall's inequality,
    \begin{equation}\label{eq:phi-gronwall}
      \norm{\Phi_\e(t)}_{L^2_r}
      \le \norm{\Phi_0}_{L^2_r}\,e^{C(\e)t}
      \quad\text{for all } t\ge 0.
    \end{equation}
  \end{enumerate}
  
  \medskip\noindent
  \textbf{The open condition.}\enspace
  If $C(\e)$ admits an $\e$-independent upper bound
  $C_0 = C_0(M, G_0)$, independent of $\e > 0$, then the Leray--Hopf
  weak solution of Theorem~\ref{thm:limit} is globally smooth.
  This is Open Problem~\ref{op:gronwall}.
\end{theorem}

The proof of Theorem~\ref{thm:phi} is given in full in the next section.

%====================================================================
\section{The Phi-Energy Estimate in Full}\label{sec:phienergy}
%====================================================================

\subsection{Taking the \texorpdfstring{$L^2_r$}{L2r} inner product}

Take the $L^2_r$ inner product of \eqref{eq:NSe-Phi} with $\Phi_\e$,
i.e., multiply by $\Phi_\e$ and integrate against $r\,dr\,dz$:
\begin{align}
  \half\ddt\norm{\Phi_\e}_{L^2_r}^2
  + \underbrace{\int\!\!\left(u_\e^r\partial_r\Phi_\e
    + u_\e^z\partial_z\Phi_\e
    + \frac{2u_\e^r\Phi_\e}{r}\right)\!\Phi_\e\,r\,dr\,dz}_{=:\,T_{\rm adv}}
  &\nonumber\\
  = \nu\underbrace{\int\!\!\Lc_4\Phi_\e\cdot\Phi_\e\,r\,dr\,dz}_{=:\,T_{\rm visc}}
  &- \e\norm{\Phi_\e}_{\dot H^{2.6}_r}^2.
  \label{eq:Phienergy}
\end{align}

\subsection{The advection term \texorpdfstring{$T_{\rm adv}$}{T-adv}}

\begin{lemma}\label{lem:Tadv}
  $T_{\rm adv} = 2\displaystyle\int\! u_\e^r\,\Phi_\e^2\,dr\,dz$.
\end{lemma}
\begin{proof}
  Write $T_{\rm adv} = \frac{1}{2}\int u_\e\cdot\nabla(\Phi_\e^2)
  \cdot r\,dr\,dz + 2\int u_\e^r\Phi_\e^2\,dr\,dz$.
  Integrate by parts:
  \[
    \frac{1}{2}\int u_\e^r\partial_r(\Phi_\e^2)\cdot r\,dr\,dz
    = -\frac{1}{2}\int\Phi_\e^2\,\partial_r(r u_\e^r)\,dr\,dz,
  \]
  \[
    \frac{1}{2}\int u_\e^z\partial_z(\Phi_\e^2)\cdot r\,dr\,dz
    = -\frac{1}{2}\int\Phi_\e^2\,\partial_z(r u_\e^z)\,dr\,dz.
  \]
  By axisymmetric incompressibility \eqref{eq:incomp},
  $\partial_r(r u_\e^r) + \partial_z(r u_\e^z) = 0$.  Hence the two
  integration-by-parts terms sum to zero, leaving
  $T_{\rm adv} = 2\int u_\e^r\Phi_\e^2\,dr\,dz$.
\end{proof}

\subsection{Bounding \texorpdfstring{$T_{\rm adv}$}{T-adv}}

Factor out $u_\e^r/r$:
\begin{equation}\label{eq:Tadvbound}
  |T_{\rm adv}|
  = 2\left|\int\frac{u_\e^r}{r}\,\Phi_\e^2\,r\,dr\,dz\right|
  \le 2\norm{\frac{u_\e^r}{r}}_{L^\infty}\norm{\Phi_\e}_{L^2_r}^2.
\end{equation}
For smooth axisymmetric fields, the axis regularity condition
$u^r|_{r=0} = 0$ together with $u^r = O(r)$ implies $u^r/r$ is bounded
and smooth.  For the $Q_1$-augmented system, Theorem~\ref{thm:lions}
ensures $u_\e \in C^\infty$ for each $\e > 0$, so
\begin{equation}\label{eq:Ce}
  C(\e) := 2\sup_{t\in[0,T]}\norm{\frac{u_\e^r(t)}{r}}_{L^\infty} < \infty
  \quad\text{for each fixed }\e > 0.
\end{equation}
Equation \eqref{eq:Tadvbound} becomes $|T_{\rm adv}| \le C(\e)\norm{\Phi_\e}_{L^2_r}^2$.

\subsection{The viscous term \texorpdfstring{$T_{\rm visc}$}{T-visc}}

\begin{lemma}\label{lem:Tvisc}
  $T_{\rm visc}
  = -\norm{\nabla\Phi_\e}_{L^2_r}^2
  - \displaystyle\int_{\R}\!\Phi_\e(0,z)^2\,dz \le 0$.
\end{lemma}
\begin{proof}
  Recall $\Lc_4 = \partial_{rr} + (3/r)\partial_r + \partial_{zz}$.
  We compute each piece against the weight $r\,dr\,dz$.

  \medskip\noindent
  \emph{The $\partial_{rr}$ piece.}
  Integrating by parts in $r$,
  \[
    \int\partial_{rr}\Phi_\e\cdot\Phi_\e\cdot r\,dr\,dz
    = -\int(\partial_r\Phi_\e)^2\,r\,dr\,dz
      - \int\partial_r\Phi_\e\cdot\Phi_\e\,dr\,dz,
  \]
  the boundary contribution
  $[\,r\,\partial_r\Phi_\e\,\Phi_\e\,]_{r=0}^{r=\infty}$ vanishing (the
  weight $r$ kills the axis endpoint, decay kills $r=\infty$).

  \medskip\noindent
  \emph{The $(3/r)\partial_r$ piece.}
  \[
    3\int\frac{1}{r}\partial_r\Phi_\e\cdot\Phi_\e\cdot r\,dr\,dz
    = 3\int\partial_r\Phi_\e\cdot\Phi_\e\,dr\,dz.
  \]

  \medskip\noindent
  \emph{The $\partial_{zz}$ piece.}
  \[
    \int\partial_{zz}\Phi_\e\cdot\Phi_\e\cdot r\,dr\,dz
    = -\int(\partial_z\Phi_\e)^2\,r\,dr\,dz.
  \]

  Summing the three pieces,
  \[
    T_{\rm visc}
    = -\norm{\nabla\Phi_\e}_{L^2_r}^2
      + 2\int\partial_r\Phi_\e\cdot\Phi_\e\,dr\,dz .
  \]
  The remaining term is a genuine axis boundary term and does
  \emph{not} vanish in general:
  \[
    2\int\partial_r\Phi_\e\cdot\Phi_\e\,dr\,dz
    = \int_{\R}\!\bigl[\Phi_\e^2\bigr]_{r=0}^{r=\infty}\,dz
    = -\int_{\R}\!\Phi_\e(0,z)^2\,dz,
  \]
  since $\Phi_\e\to 0$ as $r\to\infty$ while
  $\Phi_\e(0,z) = \partial_r u^\theta_\e|_{r=0}$ need not vanish.  Hence
  \[
    T_{\rm visc} = -\norm{\nabla\Phi_\e}_{L^2_r}^2
    - \int_{\R}\!\Phi_\e(0,z)^2\,dz \le 0,
  \]
  both terms being non-positive.  The axis term is \emph{favorable}---it
  supplies additional dissipation---so the estimate only improves.
  \qedhere
\end{proof}

\subsection{Assembling the Phi-energy estimate}

Substituting Lemmas~\ref{lem:Tadv} and~\ref{lem:Tvisc} into
\eqref{eq:Phienergy}:
\begin{equation}\label{eq:phifull}
  \half\ddt\norm{\Phi_\e}_{L^2_r}^2
  + \nu\norm{\nabla\Phi_\e}_{L^2_r}^2
  + \nu\int_{\R}\!\Phi_\e(0,z)^2\,dz
  + \e\norm{\Phi_\e}_{\dot H^{2.6}_r}^2
  \le C(\e)\norm{\Phi_\e}_{L^2_r}^2.
\end{equation}
Gronwall's lemma applied to
$\ddt\norm{\Phi_\e}_{L^2_r}^2 \le 2C(\e)\norm{\Phi_\e}_{L^2_r}^2$ gives
\[
  \norm{\Phi_\e(t)}_{L^2_r}^2
  \le \norm{\Phi_0}_{L^2_r}^2\,e^{2C(\e)t},
\]
completing the proof of Theorem~\ref{thm:phi}.

\subsection{Nature of the constant \texorpdfstring{$C(\e)$}{C(epsilon)}}

The constant $C(\e) = 2\sup_t\norm{u_\e^r/r}_{L^\infty}$ depends on the
$H^s$ regularity of $u_\e$ for $s \ge 2$, which is controlled by the
Lions theorem for each fixed $\e > 0$ but may grow as $\e\to 0$.
Controlling $\norm{u^r/r}_{L^\infty}$ uniformly in $\e$ is precisely the
open condition: it is related to Prodi--Serrin type criteria for
Navier--Stokes regularity \cite{Serrin1962}.  The connection is made
explicit in Open Problem~\ref{op:gronwall}.

%====================================================================
\section{Global Regularity of the \texorpdfstring{$\e$}{epsilon}-System}
\label{sec:global}
%====================================================================

\begin{theorem}[Global smooth solutions for $\mathrm{NS}_\e$]%
  \label{thm:uniform}
  Let $\beta = 0.6$ and $\e > 0$.  For axisymmetric-with-swirl initial
  data $u_0 \in H^\infty_{\rm axi}(\R^3)$ with $\Gamma_0\in L^\infty$
  and $\Phi_0 \in L^2_r$, the $Q_1$-augmented system \eqref{eq:NSe}
  admits a unique global smooth solution
  $u_\e\in C^\infty(\R^3\times[0,\infty))$ satisfying:
  \begin{itemize}
    \item $\norm{u_\e(t)}_{L^2} \le M$ \quad (uniform in $\e$ and $t$),
    \item $\displaystyle\int_0^T\norm{\nabla u_\e}_{L^2}^2\,dt
      \le M^2/(2\nu)$ \quad (uniform in $\e$),
    \item $\norm{\Gamma_\e(t)}_{L^\infty} \le G_0$ \quad (uniform in $\e$ and $t$),
    \item $\norm{\Phi_\e(t)}_{L^2_r} \le \norm{\Phi_0}_{L^2_r}
      e^{C(\e)t}$ \quad (for each fixed $\e > 0$).
  \end{itemize}
\end{theorem}
\begin{proof}
  Global existence and smoothness follow from Theorem~\ref{thm:lions}
  ($\beta > \frac{1}{2}$).  The uniform $L^2$ bound is
  Proposition~\ref{prop:L2}.  The $\Gamma$-bound is Lemma~\ref{lem:Gamma}.
  The $\Phi$-bound is Theorem~\ref{thm:phi}(iii).

  Higher Sobolev bounds follow by a standard elliptic bootstrap:
  given $H^k$ control on $(\Omega_\e, \Gamma_\e, \Phi_\e)$, the
  equations \eqref{eq:NSe-Omega}--\eqref{eq:NSe-Phi} provide $H^{k+1}$
  bounds by differentiating the equations and applying the same energy
  method, with constants depending on $\e$ through $\norm{u_\e}_{H^k}$.
  The axisymmetric-with-swirl structure is preserved: the axis conditions
  $u^r = u^\theta = 0$ at $r=0$ are consequences of the symmetry class
  and the smoothness of $u_\e$, and propagate to all times.
\end{proof}

%====================================================================
\section{Compactness and the Limit \texorpdfstring{$\e\to 0$}{epsilon to 0}}
\label{sec:compact}
%====================================================================

\subsection{Setup}

Fix $T > 0$.  The family $\{u_\e\}_{0<\e\le 1}$ satisfies:
\begin{align}
  &\text{(U1)}\quad \sup_{t\in[0,T]}\norm{u_\e(t)}_{L^2} \le M,
    \label{U1}\\
  &\text{(U2)}\quad \int_0^T\norm{\nabla u_\e}_{L^2}^2\,dt
    \le M^2/(2\nu),
    \label{U2}
\end{align}
both uniform in $\e$.

\subsection{Aubin--Lions compactness}

\begin{lemma}[Aubin--Lions]\label{lem:aubin}
  The family $\{u_\e\}$ is precompact in $L^2([0,T];L^2_{\rm loc}(\R^3))$.
\end{lemma}
\begin{proof}
  From (U1)--(U2), $\{u_\e\}$ is bounded in
  $L^2([0,T];H^1_{\rm axi}(\R^3))$.

  For the time derivative: testing \eqref{eq:NSe} against
  $\phi\in H^1_{\rm axi}$ with $\norm{\phi}_{H^1}\le 1$,
  \begin{align*}
    |\inner{\partial_t u_\e}{\phi}|
    &\le |\inner{(u_\e\cdot\nabla)u_\e}{\phi}|
      + \nu|\inner{\nabla u_\e}{\nabla\phi}|
      + \e|\inner{(-\Delta)^{1.3}u_\e}{\phi}|\\
    &\le \norm{u_\e}_{L^4}^2\norm{\phi}_{H^1}
      + \nu\norm{\nabla u_\e}_{L^2}\norm{\phi}_{H^1}
      + \e\norm{u_\e}_{\dot H^{2.6}}\norm{\phi}_{\dot H^{2.6}}.
  \end{align*}
  Integrating in $t$ and using (U1)--(U2) and Sobolev embedding
  $H^1\hookrightarrow L^4$ in $\R^3$:  $\{\partial_t u_\e\}$ is
  bounded in $L^{4/3}([0,T];H^{-1}_{\rm axi})$.

  By the Aubin--Lions lemma \cite{Aubin1963,Lions1969}, a bounded set in
  $L^2([0,T];H^1)\cap W^{1,4/3}([0,T];H^{-1})$ is precompact in
  $L^2([0,T];L^2_{\rm loc})$.
\end{proof}

\subsection{The limit is a Leray--Hopf weak solution}

\begin{theorem}[Weak limit]\label{thm:limit}
  There exists a subsequence $\e_k\to 0$ and a divergence-free
  axisymmetric-with-swirl vector field $u\in L^\infty([0,T];L^2)\cap
  L^2([0,T];H^1)$ such that:
  \begin{enumerate}
  \item[\textup{(i)}]
    $u_{\e_k}\to u$ strongly in $L^2([0,T];L^2_{\rm loc}(\R^3))$,
  \item[\textup{(ii)}]
    $u$ is a Leray--Hopf weak solution of the classical system
    \eqref{eq:NS} with initial data $u_0$.
  \end{enumerate}
\end{theorem}
\begin{proof}
  \textbf{Step 1: Extraction.}
  By Lemma~\ref{lem:aubin}, extract a subsequence
  $u_{\e_k}\to u$ strongly in $L^2([0,T];L^2_{\rm loc})$ and weakly
  in $L^2([0,T];H^1)$.

  \textbf{Step 2: Limit of the nonlinear term.}
  For $\phi\in C_0^\infty(\R^3\times[0,T])$:
  \begin{align*}
    \int_0^T\inner{(u_{\e_k}\cdot\nabla)u_{\e_k}}{\phi}\,dt
    -\int_0^T\inner{(u\cdot\nabla)u}{\phi}\,dt
    &= \int_0^T\inner{(u_{\e_k}-u)\cdot\nabla u_{\e_k}}{\phi}\,dt\\
    &\quad + \int_0^T\inner{u\cdot\nabla(u_{\e_k}-u)}{\phi}\,dt.
  \end{align*}
  The first term: $\norm{(u_{\e_k}-u)\cdot\nabla u_{\e_k}}_{L^1_{t,x}}
  \le \norm{u_{\e_k}-u}_{L^2_t L^4_x}\norm{\nabla u_{\e_k}}_{L^2_{t,x}}\to 0$
  by strong convergence in $L^2_t L^4_x$ (from strong $L^2_{\rm loc}$
  and the Sobolev embedding) and (U2).
  The second term goes to zero by weak convergence of $u_{\e_k}$ in
  $L^2_t H^1$.

  \textbf{Step 3: Limit of the $Q_1$ term.}
  For test function $\phi\in C_0^\infty$, the Cauchy--Schwarz inequality and
  \eqref{eq:H1int} give:
  \begin{align*}
    \left|\int_0^T\e_k\inner{(-\Delta)^{1.3}u_{\e_k}}{\phi}\,dt\right|
    &\le \e_k^{1/2}
      \left(\e_k\int_0^T\norm{u_{\e_k}}_{\dot H^{2.6}}^2\,dt\right)^{1/2}
      \left(\int_0^T\norm{\phi}_{\dot H^{2.6}}^2\,dt\right)^{1/2}\\
    &\le \e_k^{1/2}\cdot\frac{M}{\sqrt{2}}
      \cdot\norm{\phi}_{L^2_t \dot H^{2.6}} \to 0.
  \end{align*}

  \textbf{Step 4: Limit of the linear terms.}
  The pressure and viscous terms pass to the limit by weak convergence.

  \textbf{Step 5: Energy inequality.}
  The Leray--Hopf energy inequality for $u$ follows from lower
  semicontinuity of the $L^2$ and $H^1$ norms:
  \[
    \norm{u(t)}_{L^2}^2 + 2\nu\int_0^t\norm{\nabla u}_{L^2}^2\,ds
    \le \norm{u_0}_{L^2}^2.
  \]

  \textbf{Step 6: Axisymmetric structure.}
  The strong $L^2_{\rm loc}$ limit of axisymmetric-with-swirl fields is
  axisymmetric-with-swirl.
\end{proof}

\begin{remark}
  Theorem~\ref{thm:limit} produces a Leray--Hopf weak solution for
  axisymmetric-with-swirl data.  In the without-swirl case ($u^\theta\equiv 0$),
  the regularity of such solutions is known
  \cite{Ladyzhenskaya,UY1968}.  In the with-swirl case, regularity depends
  on whether $C(\e)$ in \eqref{eq:phi-energy} is $\e$-independent.
\end{remark}

%====================================================================
\section{Open Problems}\label{sec:open}
%====================================================================

\begin{enumerate}

\item[\textbf{OP1.}]\label{op:limits}
  \textbf{(Surjectivity of $Q_1$-approximation.)}
  Does every Leray--Hopf weak solution with axisymmetric-with-swirl data
  in $H^1$ arise as the strong $L^2$-limit of a sequence of
  $\mathrm{NS}_\e$ smooth solutions?  This would confirm that the class
  studied in Theorem~\ref{thm:limit} is exhaustive.

\item[\textbf{OP2.}]\label{op:gronwall}
  \textbf{(Uniform Gronwall bound --- the precise gap.)}
  The constant $C(\e)$ in \eqref{eq:phi-energy} equals
  $2\sup_t\norm{u_\e^r(t)/r}_{L^\infty}$.  The open condition is:
  \[
    \sup_{\e\in(0,1]}\int_0^T\norm{\frac{u_\e^r(t)}{r}}_{L^\infty}\,dt
    \le C_0(T, M, G_0) < \infty.
  \]
  This is a Prodi--Serrin type criterion \cite{Serrin1962} for the family
  $\{u_\e\}$: it asserts that the radial component is controlled near
  the axis uniformly in the augmentation parameter.  For smooth axisymmetric
  fields, $u^r/r = -\partial_z u^z/2 + O(r^2)$ at the axis (by
  incompressibility and Taylor expansion), so the condition is equivalent
  to a uniform bound on $\partial_z u^z_\e$ at $r=0$.

\end{enumerate}

%====================================================================
\section{Conclusion}\label{sec:conclusion}
%====================================================================

We have presented a complete treatment of the $Q_1$-augmented
axisymmetric-with-swirl Navier--Stokes system, with the following
rigorously established results:

\begin{enumerate}
  \item The substitution $\Phi = u^\theta/r$ with $\Gamma = r^2\Phi$
    yields the identity $r^{-4}\partial_z(\Gamma^2) = \partial_z(\Phi^2)$
    (Lemma~\ref{lem:cancel}), removing the explicit $1/r^4$ prefactor from
    the swirl-vorticity source and relocating the axis difficulty into the
    advection coefficient $u^r/r$.
    
  \item The Phi-energy estimate (Theorem~\ref{thm:phi}) closes for each
    fixed $\e > 0$ with constant $C(\e) = 2\sup_t\norm{u_\e^r/r}_{L^\infty}$.
    
  \item The augmented family is uniformly bounded in $L^\infty_t L^2$
    and $L^2_t H^1$, and the $\Gamma$-bound is uniform in $\e$
    (Lemma~\ref{lem:Gamma}).
    
  \item Aubin--Lions compactness (Lemma~\ref{lem:aubin}) yields
    a subsequential limit that is a Leray--Hopf weak solution of the
    classical system (Theorem~\ref{thm:limit}).
    
  \item The $Q_1$ augmentation term vanishes in the limit via the
    energy bound $\e_k\int\norm{u_{\e_k}}_{\dot H^{2.6}}^2 \le M^2/2$
    (proved in Theorem~\ref{thm:limit}, Step~3).
\end{enumerate}

\noindent
The axisymmetric-with-swirl regularity problem has been reduced to a
single open condition: $\e$-independence of the Gronwall constant
$C(\e) = 2\sup_t\norm{u_\e^r(t)/r}_{L^\infty}$ (Open Problem~\ref{op:gronwall}).

%====================================================================
\section*{Acknowledgments}
%====================================================================

The author thanks the mathematical community for the foundational
contributions of J.~Leray, J.-L.~Lions, J.~Aubin, O.~A.~Ladyzhenskaya,
M.~R.~Ukhovskii, V.~I.~Yudovich, and J.~Serrin.

%====================================================================
\begin{thebibliography}{99}
%====================================================================

\bibitem{Aubin1963}
J.~Aubin (1963).
Un th\'{e}or\`{e}me de compacit\'{e}.
\textit{C.\ R.\ Acad.\ Sci.\ Paris} \textbf{256}, 5042--5044.

\bibitem{CaoTiti2008}
C.~Cao and E.~S.~Titi (2008).
Regularity criteria for the three-dimensional Navier--Stokes equations.
\textit{Indiana Univ.\ Math.\ J.} \textbf{57}(6), 2643--2661.

\bibitem{CF1988}
P.~Constantin and C.~Foias (1988).
\textit{Navier--Stokes Equations}.
University of Chicago Press.

\bibitem{FeffermanClay}
C.~L.~Fefferman (2000).
Existence and smoothness of the Navier--Stokes equation.
In \textit{The Millennium Prize Problems}, Clay Mathematics Institute.

\bibitem{HouLi2008}
T.~Y.~Hou and C.~Li (2008).
Dynamic stability of the three-dimensional axisymmetric Navier--Stokes
equations with swirl.
\textit{Comm.\ Pure Appl.\ Math.} \textbf{61}(5), 661--697.

\bibitem{Ladyzhenskaya}
O.~A.~Ladyzhenskaya (1969).
\textit{The Mathematical Theory of Viscous Incompressible Flow}.
Gordon and Breach.

\bibitem{Leray1934}
J.~Leray (1934).
Sur le mouvement d'un liquide visqueux emplissant l'espace.
\textit{Acta Math.} \textbf{63}, 193--248.

\bibitem{LeiZhang}
Z.~Lei and Q.~S.~Zhang (2011).
Criticality of the axially symmetric Navier--Stokes equations.
\textit{Pacific J.\ Math.} \textbf{249}(1), 179--197.

\bibitem{Lions1969}
J.-L.~Lions (1969).
\textit{Quelques m\'{e}thodes de r\'{e}solution des probl\`{e}mes aux
limites non lin\'{e}aires}.
Dunod, Paris.

\bibitem{MajdaBertozzi}
A.~J.~Majda and A.~L.~Bertozzi (2002).
\textit{Vorticity and Incompressible Flow}.
Cambridge University Press.

\bibitem{Serrin1962}
J.~Serrin (1962).
On the interior regularity of weak solutions of the Navier--Stokes
equations.
\textit{Arch.\ Rat.\ Mech.\ Anal.} \textbf{9}, 187--195.

\bibitem{Stinga2010}
P.~R.~Stinga and J.~L.~Torrea (2010).
Extension problem and Harnack's inequality for some fractional operators.
\textit{Comm.\ Partial Differential Equations}
\textbf{35}(11), 2092--2122.

\bibitem{UY1968}
M.~R.~Ukhovskii and V.~I.~Yudovich (1968).
Axially symmetric flows of ideal and viscous fluids filling the whole space.
\textit{J.\ Appl.\ Math.\ Mech.} \textbf{32}, 52--61.

\end{thebibliography}

\end{document}
